\clearpage
\fancyhead[L]{\small COMET.}
\hypertarget{native-456-comet}{}
\begin{center}
{\large COMET.}
\end{center}

\begin{center}
\rule{6cm}{0.4pt}\\
Göttingische gelehrte Anzeigen. Stück 57. Seite 561..565. 1826. April 10.\\
\rule{6cm}{0.4pt}
\end{center}

Von der wichtigen Entdeckung, die uns einen Cometen von nur dreijähriger
Umlaufszeit kennen gelehrt hat (nach dem Astronomen, welcher sich um seine
Theorie so sehr verdient gemacht hat, der
\SourceLeaf{457}\textsc{Encke}sche Comet genannt), hatten unsere Blätter im Jahre 1819 St. 28 und 83.
[S. 417 u. 420 d. B.] die erste öffentliche Nachricht gegeben, die sich mit
dem Wunsche schloss, dass \emph{diese höchst merkwürdige Entdeckung die
Freunde der Astronomie zu verdoppeltem Eifer in Aufsuchung von Cometen
anfeuern möchte, da alles vermuthen lasse, dass jene nur der Anfang einer
unermesslichen nach und nach reifenden Erndte sein werde.} In welchem Grade
jener \emph{Wunsch} in Erfüllung gegangen ist, beweisen 14 seitdem
beobachtete und berechnete neue Cometen: mit einem hohen Grade von
Wahrscheinlichkeit können wir jetzt auch die erste Bestätigung der
\emph{Vermuthung} anzeigen. Der Herr Hauptmann \textsc{Biela} in Josephsstadt in
Böhmen entdeckte am 28. Februar d.\,J. im Sternbilde des Widders einen
kleinen, dem blossen Auge unsichtbaren Cometen, welcher auf die von ihm
gemachte Anzeige auch an mehrern andern Orten, auch bisher zweimal auf
hiesiger Sternwarte durch Hrn. Prof. \textsc{Harding}, beobachtet worden ist;
denselben Cometen hat auch einige Tage später Herr \textsc{Gambard} in Marseille für
sich entdeckt. Bei der grossen Menge neuer Cometen, die in den letzten Jahren
aufgefunden sind, würden wir des gegenwärtigen in diesen Blättern nicht
besonders erwähnen, wenn nicht schon die ersten von mehrern verschiedenen
Berechnern, auch auf verschiedenen Beobachtungen gegründeten Rechnungen über
seine Bahn ein höchst merkwürdiges Resultat gegeben hätten. Hr. Hofr. \textsc{Gauss}
hat die parabolischen von fünf Astronomen, dem Hrn. Hauptmann von \textsc{Biela}
selbst, dem Hrn. \textsc{Clausen} in Altona, Hrn. Prof. \textsc{Encke} in Berlin, Hrn.
Dr. \textsc{Olbers} in Bremen und Hrn. Prof. \textsc{Schwerdt} in Speyer mitgetheilt erhalten,
die wir hier zusammenstellen:

\begin{center}
\begingroup\small\setlength{\tabcolsep}{5pt}
\begin{tabular}{rlll}
\multicolumn{4}{l}{Durchgang durch die Sonnennähe} \\
1826. März & 15.45653 & v. B. & \\
 & 15.75422 & Cl. & \\
 & 18.2432 & Göttinger Z. & E. \\
 & 18.33773 & --- & O. \\
 & 17.620 & Sch. & \\
\end{tabular}
\par\medskip
\begin{tabular}{ll}
\multicolumn{2}{l}{Länge der Sonnennähe} \\
$96^{\circ}\quad27'\quad33$ & v. B. \\
$95\quad48\quad15$ & Cl. \\
$107\quad55\quad31$ & E. \\
$112\quad39\quad24$ & O \\
$100\quad51\quad57$ & Sch. \\
\end{tabular}
\par\medskip
\begin{tabular}{ll}
\multicolumn{2}{l}{Aufsteigender Knoten} \\
$245^{\circ}\quad57'\quad0''$ & v. B. \\
$245\quad16\quad40$ & Cl. \\
$248\quad41\quad59$ & E. \\
$250\quad44\quad21$ & O. \\
$246\quad34\quad52$ & Sch. \\
\end{tabular}
\par\medskip
\begin{tabular}{ll}
\multicolumn{2}{l}{Neigung der Bahn} \\
$15^{\circ}\quad28'\quad23''$ & v. B. \\
$15\quad45\quad19$ & Cl. \\
$13\quad12\quad37$ & E. \\
$12\quad18\quad31$ & O. \\
$15\quad7\quad5$ & Sch. \\
\end{tabular}
\par\medskip
\begin{tabular}{ll}
\multicolumn{2}{l}{Logarithm des kleinsten Abstandes} \\
0.00506 & v. B. \\
0.0114369 & Cl. \\
9.96474 & E. \\
9.94460 & O. \\
9.99565 & Sch. \\
\end{tabular}
\endgroup
\end{center}
\noindent Bewegung rechtläufig.

Die Unterschiede dieser Resultate sind nicht bedeutend, in so fern man die
Dürftigkeit der ihnen zum Grunde liegenden Data erwägt, und werden sich bald
ausgleichen lassen.

Sogleich höchst auffallend ist nun die grosse Aehnlichkeit, welche diese
Elemente durchgehends mit denen des zweiten vom Jahre 1805 haben. Wenn schon
diese allein eine ungemein grosse Wahrscheinlichkeit für die Identität beider
Cometen begründet, so wird diese Wahrscheinlichkeit noch ausserordentlich
durch die Umstände verstärkt, die den Cometen von 1805 betreffen. Hr. Hofr.
\textsc{Gauss}, welcher über dessen Erscheinung eine sehr ausgedehnte Untersuchung
ausgeführt hatte, fand, dass die
\SourceLeaf{458}Beobachtungen sehr auffallend auf eine kurze Umlaufszeit hinwiesen, obgleich sie nicht
zureichten, solche aus dieser Erscheinung allein mit einiger Genauigkeit festzusetzen:
am besten liessen sich die Beobachtungen mit einer Umlaufszeit von $4\frac{3}{4}$ Jahren
vereinigen; aber eine ein Jahr kürzere oder einige Jahre längere Umlaufszeit, hätte sich,
jenen Untersuchungen zufolge, gleichfalls noch sehr gut und viel besser als die
parabolische Bewegung mit den Beobachtungen in Uebereinstimmung bringen lassen.
Ausserdem hatte sich bei den Elementen des Cometen von 1805 eine grosse Aehnlichkeit
mit denen des Cometen von 1772 gezeigt, und auf die Vermuthung geführt, dass beide
identisch sein möchten; inzwischen war der Unterschied doch noch so beträchtlich, dass
einige Astronomen die Identität für unwahrscheinlich oder selbst die Verschiedenheit für
gewiss hielten. Hr. Hofr. \textsc{Gauss} zeigte aber schon damals, dass dies Urtheil nicht
hinlänglich begründet, und dass die Möglichkeit einer vollkommenen Erklärung jenes
Unterschieds keineswegs ausgeschlossen sei, und wenn gleich er in der Zwischenzeit
nicht dazu gekommen ist, die weitläuftigen zur Entscheidung nöthigen Rechnungen zu
unternehmen, so ist ihm doch fortwährend die Identität beider Cometen wahrscheinlich
geblieben. Unter diesen Umständen ist es nun höchst merkwürdig, dass die Zwischenzeiten
einerseits zwischen den Durchgängen der Cometen von 1772 und 1805 durch ihre
Sonnennähe, und andererseits zwischen den Durchgängen der Cometen von 1805 und
1826 nahe im Verhältnisse der Zahlen 5 und 3 stehen.

Mit hoher Wahrscheinlichkeit dürfen wir daher vermuthen, dass alle drei Cometen
identisch sind, dass dieser Comet von 1772 bis 1826 zusammen acht Umläufe gemacht,
jeden im Durchschnitt von 6 Jahren und 9 Monaten, und dass die Unterschiede in den
Elementen von 1772 und 1805, durch welche einige Astronomen an der vermutheten
Identität irre geworden waren, hauptsächlich eine Folge der Störungen durch den
Jupiter gewesen sind, dem der Comet in der Zwischenzeit sehr nahe gekommen ist. Ohne
diesen letzten Umstand würde es sonst für den Augenblick noch eben so wahrscheinlich
sein, dass der Comet in der Zwischenzeit 10 und 6, also zusammen 16 jeden von
$3\frac{3}{8}$ Jahren gemacht hätte.

Es ist sehr zu wünschen, dass der von Hrn. v. \textsc{Biela} entdeckte Comet, dessen
Beobachtungen in der Mitte des März durch den Mondschein unterbrochen wurden, nach
Aufhören desselben noch wiederholt irgendwo gut beobachtet werden möge. Sollte dies
aber auch nicht gelingen, so werden selbst die schon gewonnenen Beobachtungen, wenn
sie erst einer angemessenen Behandlung unterworfen werden, schon zureichen, über die
sich so unwiderstehlich aufdringende Vermuthung der Identität dieser drei Cometen zu
entscheiden.
